\section{Introduction}
\label{sec:intro}

Transformer MLP layers are the primary mechanism by which language models store and transform factual knowledge~\citep{geva2021transformer,meng2022locating,dai2022knowledge}.
A growing body of work shows that individual MLP neurons promote specific tokens in vocabulary space~\citep{geva2022transformer}, that MLP layers implement vector arithmetic for factual recall~\citep{merullo2023language}, and that intermediate representations are organized as linear directions corresponding to concepts~\citep{park2024linear}.
These findings suggest an appealing hypothesis: perhaps MLP computation can be described as a \emph{simple program} that checks which token directions are present in the input and produces a shallow combination of output token directions.

If true, such a description would make transformer internals far more interpretable.
Rather than treating each MLP layer as an opaque $\dmodel \to \dmodel$ nonlinear function with $\dmlp$ hidden units, we could describe it as: ``given that the input contains directions for tokens $\{t_1, t_2, \ldots\}$, produce an update that promotes tokens $\{t'_1, t'_2, \ldots\}$.''
This would enable targeted model editing, debugging, and verification at a level of abstraction that connects directly to the model's vocabulary.

{\bf What is missing?}
Despite the evidence above, no prior work has systematically quantified how much of the \emph{full} MLP input-to-output mapping can be captured by a shallow program operating over token directions.
\citet{geva2021transformer} show that MLP keys detect input patterns and \citet{geva2022transformer} show that values promote concepts, but these analyses examine individual neurons in isolation.
\citet{merullo2023language} demonstrate vector arithmetic for specific relational tasks, but do not measure how much of the aggregate MLP behavior this explains.
We lack a direct answer to the question: \emph{what fraction of MLP computation is describable as a linear program over token directions?}

We address this gap by projecting MLP inputs and outputs into vocabulary space via the unembedding matrix $\Wu$ and fitting sparse linear models that map input token logits to output token logits.
We evaluate this ``token-direction program'' across all 12 layers of \gpttwo on 19,281 token positions from \wikitext, comparing against random-direction and full-rank linear baselines.

Our key finding is that the token-direction program hypothesis is \emph{partially supported with significant caveats}.
At layer~11, the program achieves $\Rsq=0.54$ (or $0.70$ with $k=500$ directions), but middle layers resist any linear description ($\Rsq < 0.15$).
Critically, random orthogonal directions perform comparably to token directions, revealing that the approximability comes from the MLP's inherent partial linearity---not from token directions being a privileged basis.

Our contributions are as follows:
\begin{itemize}[leftmargin=*,itemsep=0pt,topsep=0pt]
    \item We propose a systematic framework for testing whether MLP layers can be described as shallow programs over token directions, decomposing MLP inputs and outputs via the unembedding matrix and fitting sparse linear approximations.
    \item We discover a U-shaped linearity profile across layers: early and late layers are approximately linear in any basis, while middle layers are highly nonlinear.
    \item We demonstrate via a random-direction control that token directions are not a privileged basis for describing MLP computation---the partial success of token-direction programs reflects linear structure, not vocabulary-specific organization.
    \item We provide layer-wise quantitative benchmarks ($\Rsq$, cosine similarity, top-$k$ agreement) that characterize where the ``MLP as token program'' metaphor holds and where it breaks down.
\end{itemize}
